\documentclass[11pt,a4paper]{article}

% --- Encoding ---
\usepackage[utf8]{inputenc}
\usepackage[T1]{fontenc}
\usepackage[spanish,es-noshorthands]{babel}
\usepackage{lmodern}

% --- Layout ---
\usepackage[margin=2.5cm]{geometry}
\usepackage{parskip}
\setlength{\parindent}{0pt}
\setlength{\parskip}{0.5em plus 0.2em minus 0.1em}

% --- Tables and graphics ---
\usepackage{booktabs}
\usepackage{array}
\usepackage{enumitem}
\usepackage{xcolor}
\usepackage{graphicx}
\usepackage{tikz}
\usetikzlibrary{positioning,shapes.geometric,arrows.meta,fit,backgrounds,calc,decorations.pathreplacing,trees}
\usepackage{caption}
\usepackage{subcaption}
\usepackage{amsmath,amssymb}

% --- Soporte para caracteres Unicode usados en formulas en texto plano ---
\usepackage{newunicodechar}
\newunicodechar{→}{\ensuremath{\rightarrow}}
\newunicodechar{≥}{\ensuremath{\geq}}
\newunicodechar{≤}{\ensuremath{\leq}}
\newunicodechar{×}{\ensuremath{\times}}
\newunicodechar{÷}{\ensuremath{\div}}
\newunicodechar{≈}{\ensuremath{\approx}}
\newunicodechar{⌈}{\ensuremath{\lceil}}
\newunicodechar{⌉}{\ensuremath{\rceil}}
\newunicodechar{✓}{\ensuremath{\checkmark}}

% --- Code listings ---
\usepackage{listings}
\definecolor{codegray}{rgb}{0.96,0.96,0.96}
\definecolor{codeborder}{rgb}{0.85,0.85,0.85}
\definecolor{codecomment}{rgb}{0.30,0.55,0.30}
\definecolor{codekeyword}{rgb}{0.10,0.30,0.70}
\definecolor{codestring}{rgb}{0.65,0.15,0.10}

\lstdefinestyle{cstyle}{
    language=C,
    backgroundcolor=\color{codegray},
    basicstyle=\ttfamily\footnotesize,
    breaklines=true,
    frame=single,
    rulecolor=\color{codeborder},
    showstringspaces=false,
    keywordstyle=\color{codekeyword}\bfseries,
    commentstyle=\color{codecomment}\itshape,
    stringstyle=\color{codestring},
    aboveskip=1em,
    belowskip=1em,
}

\lstdefinestyle{shellstyle}{
    language=bash,
    backgroundcolor=\color{codegray},
    basicstyle=\ttfamily\small,
    breaklines=true,
    frame=single,
    rulecolor=\color{codeborder},
    showstringspaces=false,
    keywordstyle=\color{codekeyword}\bfseries,
    commentstyle=\color{codecomment}\itshape,
    stringstyle=\color{codestring},
    aboveskip=1em,
    belowskip=1em,
    columns=fullflexible,
    upquote=true,
}

\lstdefinestyle{textstyle}{
    backgroundcolor=\color{codegray},
    basicstyle=\ttfamily\footnotesize,
    breaklines=true,
    frame=single,
    rulecolor=\color{codeborder},
    showstringspaces=false,
    aboveskip=1em,
    belowskip=1em,
    columns=fullflexible,
}

\lstset{style=textstyle}

% --- Definition / Idea / Important boxes ---
\usepackage[most]{tcolorbox}
\newtcolorbox{idea}[1][]{
    colback=orange!5!white,
    colframe=orange!60!black,
    fonttitle=\bfseries,
    title=Idea intuitiva,
    sharp corners=southwest,
    rounded corners=northwest,
    arc=2pt,
    boxrule=0.6pt,
    left=8pt,right=8pt,top=4pt,bottom=4pt,
    #1
}
\newtcolorbox{importante}[1][]{
    colback=red!5!white,
    colframe=red!50!black,
    fonttitle=\bfseries,
    title=Importante,
    sharp corners=southwest,
    rounded corners=northwest,
    arc=2pt,
    boxrule=0.6pt,
    left=8pt,right=8pt,top=4pt,bottom=4pt,
    #1
}
\newtcolorbox{respuesta}[1][]{
    colback=green!5!white,
    colframe=green!40!black,
    fonttitle=\bfseries,
    title=Respuesta,
    sharp corners=southwest,
    rounded corners=northwest,
    arc=2pt,
    boxrule=0.6pt,
    left=8pt,right=8pt,top=4pt,bottom=4pt,
    #1
}

% --- Hyperlinks ---
\usepackage[colorlinks=true,
            linkcolor=blue!60!black,
            urlcolor=blue!50!black,
            citecolor=blue!60!black]{hyperref}

% --- Color palette ---
\definecolor{c1}{HTML}{4C72B0}
\definecolor{c2}{HTML}{DD8452}
\definecolor{c3}{HTML}{55A868}
\definecolor{c4}{HTML}{8172B2}
\definecolor{cred}{HTML}{C44E52}
\definecolor{cfree}{HTML}{8C8C8C}

% --- Title block ---
\title{
    \vspace{-2.5em}
    \textbf{Recuperatorio del Parcial 1 --- 2026} \\[0.3em]
    \large Sistemas Operativos / Sistemas Operativos y Redes
}
\author{
    Tomás Rodeghiero \\
    \small UNRC --- Universidad Nacional de Río Cuarto
}
\date{2026}

\begin{document}

\maketitle

\begin{abstract}
\noindent
Resolución completa del Recuperatorio del Parcial 1 del año 2026 de
la cátedra de Sistemas Operativos (Lic.\ en Ciencias de la
Computación, UNRC). El parcial consta de cinco ejercicios que abarcan:
(i) verdadero/falso sobre llamadas al sistema básicas (\texttt{exec},
\texttt{exit}, \texttt{fork}, \texttt{wait}, \texttt{printf}, traps);
(ii) conteo de procesos creados por una cadena de \texttt{fork()} con
un \texttt{exec()} final; (iii) análisis de un programa de
comunicación padre/hijo con \texttt{pipe()} y propuesta de corrección;
(iv) análisis de un escenario clásico de \emph{deadlock} con dos
mutex tomados en orden distinto entre tareas; (v) planificación
\emph{Round Robin} con quantum = 2 sobre tres procesos. Cada
ejercicio se desarrolla con la estructura \textbf{Enunciado $\to$
Análisis $\to$ Desarrollo $\to$ Respuesta}.
\end{abstract}

\tableofcontents
\bigskip

% =====================================================================
\section{Ejercicio 1 --- V/F sobre syscalls (2 pts)}
% =====================================================================

\textbf{Enunciado.} Determinar la verdad o falsedad de las siguientes
sentencias. Justificar las falsas en forma breve, completa y con
claridad.

\subsection*{1.A --- \texttt{exec(path, args)}}

\emph{``La llamada al sistema \texttt{exec(path, args)} crea un nuevo
proceso y carga las secciones de código y datos del programa y
comienza su ejecución.''}

\begin{respuesta}
\textbf{F}. \texttt{exec} \textbf{no crea un nuevo proceso}.
Lo que hace es \emph{reemplazar} la imagen del proceso \emph{actual}:
mismo PID, mismo descriptor en el kernel, pero todo el segmento de
código, datos y stack se reescribe con los del nuevo programa. La
ejecución continúa en el \texttt{main} del nuevo programa.

Para \emph{crear} un nuevo proceso se usa \texttt{fork()}. La
combinación clásica es \texttt{fork + exec}: el padre hace
\texttt{fork()} (creando un proceso hijo idéntico) y el hijo hace
\texttt{exec()} para correr otro programa. Así es como
\texttt{bash}, por ejemplo, lanza cualquier comando.
\end{respuesta}

\subsection*{1.B --- Inicio de la ejecución}

\emph{``Un proceso siempre comienza su ejecución por la función
\texttt{main()}.''}

\begin{respuesta}
\textbf{F}. Antes de \texttt{main()} se ejecuta el \emph{startup code}
del runtime de C (CRT). El punto de entrada real del binario es
\texttt{\_start} (suministrado por la libc), que se encarga de:

\begin{itemize}[leftmargin=*,nosep]
    \item Preparar \texttt{argc} y \texttt{argv} desde la pila inicial.
    \item Inicializar las variables globales en \texttt{.data} y
        poner en cero la sección \texttt{.bss}.
    \item Llamar a los constructores estáticos (en C++) y a las
        funciones registradas con atributo \texttt{constructor}.
    \item \emph{Recién entonces} invocar a \texttt{main()}.
    \item Al volver de \texttt{main}, llamar a las funciones de
        \texttt{atexit} y a \texttt{\_exit} con el valor de retorno.
\end{itemize}

Para el programador C, \texttt{main} \emph{parece} ser el primer
código que se ejecuta, pero el binario en sí arranca varios pasos
antes.
\end{respuesta}

\subsection*{1.C --- \texttt{printf} como syscall}

\emph{``La función \texttt{printf(fmt, ...)} es una llamada al
sistema.''}

\begin{respuesta}
\textbf{F}. \texttt{printf} es una función de la \textbf{biblioteca
estándar de C} (\texttt{libc}), \emph{no} una syscall.
Internamente:

\begin{enumerate}[leftmargin=*,nosep]
    \item Formatea los argumentos según la cadena \texttt{fmt} y los
        coloca en un buffer interno (línea o bloque, según el tipo
        de stream).
    \item Eventualmente llama a la syscall real
        \texttt{write(STDOUT\_FILENO, buffer, n)} para volcar los
        bytes al kernel.
\end{enumerate}

Por eso es posible reemplazar \texttt{printf} con otra
implementación (e.g.\ versiones \texttt{nano-printf} en sistemas
embebidos) sin cambiar el kernel: la única syscall involucrada es
\texttt{write}.
\end{respuesta}

\subsection*{1.D --- \texttt{exit(n)} y el \emph{zombie}}

\emph{``La llamada al sistema \texttt{exit(n)} finaliza el proceso y
en el descriptor del proceso queda el \texttt{exit code = n} hasta
que el proceso padre ejecute \texttt{wait(\&c)}.''}

\begin{respuesta}
\textbf{V}. Cuando un proceso ejecuta \texttt{exit(n)}:

\begin{itemize}[leftmargin=*,nosep]
    \item El kernel libera los recursos del proceso (memoria,
        descriptores abiertos, etc.).
    \item \textbf{No} libera el descriptor (\texttt{task\_struct} en
        Linux) inmediatamente: lo deja en estado \texttt{ZOMBIE}
        guardando los $8$ bits inferiores de \texttt{n} en
        \texttt{exit\_code}.
    \item Le envía la señal \texttt{SIGCHLD} al padre.
    \item Cuando el padre invoca \texttt{wait(\&c)} (o
        \texttt{waitpid}), el kernel copia el exit code en
        \texttt{c}, libera el descriptor del zombie y este desaparece
        definitivamente.
\end{itemize}

Si el padre nunca llama a \texttt{wait}, el zombie queda
indefinidamente en la tabla de procesos hasta que el padre muera y
\texttt{init}/\texttt{systemd} adopte al huérfano y lo coseche.
\end{respuesta}

\subsection*{1.E --- \texttt{fork()}}

\emph{``La llamada al sistema \texttt{fork()} copia la imagen
(código, datos y stack) en la memoria del proceso invocante previa
creación de un nuevo descriptor en la tabla de procesos en el kernel
con un nuevo PID.''}

\begin{respuesta}
\textbf{F}. Hay dos problemas con la afirmación:

\begin{enumerate}[leftmargin=*]
    \item \textbf{``en la memoria del proceso invocante''} está
        confuso: la imagen se copia (lógicamente) en la memoria
        \emph{del nuevo proceso} (el hijo), no en la del invocante
        (que sería el padre, donde ya está la imagen original). El
        hijo recibe su propio espacio de direcciones.

    \item En la práctica, los SO modernos \textbf{no copian la
        memoria entera} al hacer \texttt{fork}: usan
        \textbf{copy-on-write (COW)}. El padre y el hijo comparten
        las mismas páginas físicas marcadas como solo lectura en
        ambos. Cuando \emph{cualquiera} de los dos intenta escribir,
        el hardware genera un page fault, el kernel hace una copia
        privada de esa página para el que escribe y vuelve a la
        ejecución. Así se evita el costo de copiar todo el espacio
        de direcciones, que muchas veces es seguido por un
        \texttt{exec} que descarta todo.
\end{enumerate}

Una formulación correcta sería: ``fork crea un nuevo descriptor en
la tabla de procesos con un nuevo PID y le da al nuevo proceso un
espacio de direcciones que es una copia lógica del padre,
implementada eficientemente con copy-on-write''.
\end{respuesta}

\subsection*{1.F --- Implementación de syscalls vía \emph{trap}}

\emph{``Una llamada al sistema se implementa con una instrucción de
trap que dispara el mecanismo de manejo de interrupciones/excepciones
de la CPU.''}

\begin{respuesta}
\textbf{V}. Exactamente. La instrucción de \emph{trap} es la que
permite cruzar la barrera de modo usuario $\to$ kernel:

\begin{itemize}[leftmargin=*,nosep]
    \item En x86-64: instrucción \texttt{syscall}.
    \item En ARM (AArch64): \texttt{svc \#0}.
    \item En RISC-V: \texttt{ecall}.
\end{itemize}

Cuando la CPU encuentra la instrucción, ejecuta el mismo mecanismo
que ante una interrupción/excepción: salva el estado del proceso,
cambia a modo supervisor y salta a una dirección fija (la
\emph{trap table}) donde el kernel toma el control. El kernel
identifica qué syscall pidió el proceso (el número va en un
registro convenido) y la despacha a la rutina correspondiente.

El hecho de reusar el mismo \emph{datapath} que las
interrupciones/excepciones (mismo guardado de contexto, mismo
mecanismo de retorno) es lo que hace que las syscalls sean
\emph{rápidas}: la CPU ya tiene hardware dedicado para este flujo.
\end{respuesta}

% =====================================================================
\section{Ejercicio 2 --- Cuántos procesos crea esta cadena de forks (2 pts)}
% =====================================================================

\textbf{Enunciado.} Dado el siguiente programa, ¿cuántos procesos se
crean (además del proceso original)?

\begin{lstlisting}[style=cstyle]
#include <stdio.h>
#include <unistd.h>
int main() {
    fork();                 // (1)
    fork();                 // (2)
    if (fork() == 0) {      // (3)
        exec("myprog", NULL);
    }
    return 0;
}
\end{lstlisting}

\subsection*{Análisis}

Cada llamada \texttt{fork()} \emph{duplica} el proceso que la
ejecuta: el invocante (padre) sigue después de la llamada y queda un
nuevo proceso (hijo) que comienza a ejecutar también desde la
instrucción siguiente a \texttt{fork}. Por lo tanto, después de cada
\texttt{fork()} la cantidad de procesos en juego \textbf{se duplica}.

Importante: \texttt{exec("myprog", NULL)} \textbf{no} crea nuevos
procesos. Reemplaza la imagen del proceso que la invoca por la del
binario \texttt{myprog}, sin cambiar el PID. Para contar procesos
nuevos, \texttt{exec} no suma.

\subsection*{Conteo paso a paso}

\textbf{Antes de cualquier fork}: $1$ proceso (el original, ``P0'').

\textbf{Después de fork (1)}: cada uno de los $1$ procesos se duplica
$\to$ $1 \times 2 = 2$ procesos.

\textbf{Después de fork (2)}: $2 \times 2 = 4$ procesos.

\textbf{Después de fork (3)}: $4 \times 2 = 8$ procesos.

Notar que el \texttt{if (fork() == 0)} no impide que el fork se
ejecute: el fork sí ocurre en los $4$ procesos que llegan a esa
línea. Los hijos creados (que reciben $0$) entran al \texttt{if} y
ejecutan \texttt{exec}; los padres (que reciben un PID positivo)
salen del \texttt{if} y van directo al \texttt{return 0}.

\subsection*{Árbol de procesos}

\begin{center}
\begin{tikzpicture}[
    every node/.style={font=\ttfamily\scriptsize,
                       draw,circle,minimum size=0.55cm,inner sep=0pt},
    level distance=1.1cm,
    level 1/.style={sibling distance=8cm},
    level 2/.style={sibling distance=4cm},
    level 3/.style={sibling distance=2cm}
]
\node {P0}
    child {node {P0}
        child {node {P0}
            child {node {P0}}
            child {node {P4}}
        }
        child {node {P2}
            child {node {P2}}
            child {node {P6}}
        }
    }
    child {node {P1}
        child {node {P1}
            child {node {P1}}
            child {node {P5}}
        }
        child {node {P3}
            child {node {P3}}
            child {node {P7}}
        }
    };
\end{tikzpicture}
\end{center}

Las hojas del árbol son los $8$ procesos finales (las etiquetas a la
izquierda muestran ``herencia''; sólo los PIDs distintos importan).

\subsection*{Cálculo matemático}

\quad Cantidad total de procesos = $2^k$, donde $k$ = cantidad de
\texttt{fork} encadenados.

\quad Total = $2^3 = 8$ procesos.

\quad Procesos creados (sin contar el original) = $8 - 1 = 7$.

\begin{respuesta}
Se crean \textbf{$7$ procesos nuevos} (además del original). El total
en ejecución llega a $8$ procesos.

De esos $8$:
\begin{itemize}[leftmargin=*,nosep]
    \item $4$ son los ``hijos'' del último \texttt{fork()} (los que
        reciben $0$) y ejecutan \texttt{exec("myprog", NULL)},
        cambiando su imagen.
    \item Los otros $4$ (incluido el original) salen del \texttt{if}
        sin hacer \texttt{exec} y terminan con \texttt{return 0}.
\end{itemize}

\textbf{Importante}: \texttt{exec} no genera procesos nuevos, así que
no suma al conteo.
\end{respuesta}

% =====================================================================
\section{Ejercicio 3 --- Comunicación con \texttt{pipe()}: ¿correcto? (2 pts)}
% =====================================================================

\textbf{Enunciado.} Determinar si el siguiente programa es correcto
para que los dos procesos (padre e hijo) se comuniquen mediante pipes
y la salida sea \texttt{"OK"}. En caso de considerarlo incorrecto,
dar una versión corregida.

\begin{lstlisting}[style=cstyle]
#include <stdio.h>
#include <unistd.h>
int main() {
    int p[2]; char buffer[10];
    pipe(p);
    if (fork() > 0) {
        write(p[1], "Hi", 3);   // send message
        read(p[0], buffer, 10); // read response
        printf("%s\n", buffer);
    } else {
        read(p[0], buffer, 10); // read message
        write(p[1], "OK", 3);   // send response
    }
    return 0;
}
\end{lstlisting}

\subsection*{Análisis: ¿qué está mal?}

El programa \textbf{es incorrecto}. El problema central es que se
intenta usar \textbf{un solo pipe en ambos sentidos}, pero un pipe
de UNIX es \textbf{unidireccional}: lo que se escribe en \texttt{p[1]}
puede ser leído por cualquier proceso que tenga acceso al extremo de
lectura \texttt{p[0]}, incluido el propio escritor.

\paragraph{Traza del fallo.}

\begin{enumerate}[leftmargin=*]
    \item \textbf{Padre} hace \texttt{write(p[1], "Hi", 3)}. El pipe
        ahora contiene los bytes \texttt{H i \textbackslash 0}.
    \item \textbf{Padre} hace \texttt{read(p[0], buffer, 10)}.
    \item Si el padre lee \emph{antes} de que el hijo lo haga
        (perfectamente posible, depende del scheduling), \textbf{lee
        su propio mensaje} ``Hi'' del pipe.
    \item El padre imprime \texttt{"Hi"} en lugar del \texttt{"OK"}
        esperado.
    \item Mientras tanto, el hijo se queda bloqueado en
        \texttt{read()} esperando un mensaje que ya fue consumido por
        el padre.
\end{enumerate}

Aún suponiendo que el hijo gane la carrera y lea primero ``Hi'',
quedaría un segundo problema: el hijo escribe ``OK'' en el mismo
pipe, y entonces el padre podría leer correctamente ``OK''. Pero
\emph{esto depende del scheduling}: el programa no tiene
\textbf{determinismo}, y bajo la mayoría de scheduling es más
probable que el padre lea su propio mensaje.

\begin{importante}
\textbf{Regla general}: para comunicación bidireccional entre dos
procesos hace falta \textbf{dos pipes}, uno por sentido. Un único
pipe no alcanza porque cualquier proceso con acceso al \texttt{p[0]}
puede consumir lo que el otro acaba de escribir.
\end{importante}

\subsection*{Versión corregida}

Usar \textbf{dos pipes}: \texttt{p} para parent$\to$child y
\texttt{q} para child$\to$parent. Cada proceso cierra los extremos
que no va a usar para forzar \emph{EOF} cuando el otro termine.

\begin{lstlisting}[style=cstyle]
#include <stdio.h>
#include <unistd.h>
#include <sys/wait.h>
#include <string.h>

int main() {
    int p[2], q[2];               /* p: padre->hijo, q: hijo->padre */
    char buffer[10];

    pipe(p);
    pipe(q);

    if (fork() > 0) {
        /* PADRE: escribe en p, lee de q */
        close(p[0]);              /* no lee de p */
        close(q[1]);              /* no escribe en q */

        write(p[1], "Hi", 3);     /* envia mensaje */
        close(p[1]);              /* cierra para que el hijo vea EOF */

        read(q[0], buffer, 10);   /* lee respuesta */
        printf("%s\n", buffer);
        close(q[0]);

        wait(NULL);               /* espera al hijo (evita zombie) */
    } else {
        /* HIJO: lee de p, escribe en q */
        close(p[1]);              /* no escribe en p */
        close(q[0]);              /* no lee de q */

        read(p[0], buffer, 10);   /* lee mensaje */
        close(p[0]);

        write(q[1], "OK", 3);     /* envia respuesta */
        close(q[1]);
    }
    return 0;
}
\end{lstlisting}

\subsection*{Diagrama de la comunicación corregida}

\begin{center}
\begin{tikzpicture}[every node/.style={font=\ttfamily\small},
    proc/.style={draw,rounded corners,minimum width=2.4cm,minimum height=1.2cm,align=center}]
\node[proc,fill=c1!30] (padre) at (0,0)  {Padre};
\node[proc,fill=c2!30] (hijo)  at (7,0)  {Hijo};

\draw[->, thick, c1!60!black] (padre.north east) to[bend left=30]
    node[midway,above,font=\small]{p: ``Hi''} (hijo.north west);
\draw[->, thick, c2!60!black] (hijo.south west)  to[bend left=30]
    node[midway,below,font=\small]{q: ``OK''} (padre.south east);
\end{tikzpicture}
\end{center}

\begin{respuesta}
\textbf{El programa NO es correcto}. Usa un único pipe en ambos
sentidos, lo que provoca que el padre lea su propio mensaje
``Hi'' en vez de la respuesta ``OK''. Además, el hijo puede quedar
bloqueado indefinidamente esperando un mensaje ya consumido.

\textbf{Corrección}: usar \emph{dos} pipes, uno por sentido, y cerrar
los extremos que cada proceso no usa. La versión corregida está
arriba.
\end{respuesta}

% =====================================================================
\section{Ejercicio 4 --- Transferencias bancarias y \emph{deadlock} (2 pts)}
% =====================================================================

\textbf{Enunciado.} Un sistema bancario provee la siguiente función:

\begin{lstlisting}[style=cstyle]
void transfer(Account from, Account to, double amount)
{
    mutex lock1, lock2;
    lock1 = get_lock(from);
    lock2 = get_lock(to);
    acquire(lock1);
    acquire(lock2);
    withdraw(from, amount);   // extract money
    deposit(to, amount);      // deposit money
    release(lock2);
    release(lock1);
}
\end{lstlisting}

¿Logra prevenir deadlocks? Si no, dar un escenario de dos tareas
concurrentes invocando \texttt{transfer(\ldots)} en el que pueda
ocurrir deadlock. Suponer que existen al menos las cuentas
\texttt{accountA} y \texttt{accountB}.

\subsection*{Análisis: el orden de adquisición depende de los argumentos}

La función toma siempre los locks en el orden $(\texttt{lock1},
\texttt{lock2})$, donde $\texttt{lock1}$ corresponde a la cuenta
\texttt{from} y $\texttt{lock2}$ a la cuenta \texttt{to}.

El problema: \textbf{el orden de adquisición no es global}, depende
de cuál cuenta es \texttt{from} y cuál es \texttt{to} en cada
invocación. Eso permite que dos invocaciones concurrentes
\emph{intenten tomar los mismos dos locks en orden inverso}, lo que
es exactamente la receta clásica para un deadlock.

\subsection*{Escenario concreto de deadlock}

Sean dos transferencias concurrentes ejecutadas por dos threads:

\begin{center}
\small
\begin{tabular}{@{}lp{6cm}p{6cm}@{}}
\toprule
$t$ & \textbf{Thread T1: \texttt{transfer(A, B, 100)}} & \textbf{Thread T2: \texttt{transfer(B, A, 50)}} \\
\midrule
$t_1$ & \texttt{lock1 = get\_lock(A)}   & \texttt{lock1 = get\_lock(B)} \\
$t_2$ & \texttt{lock2 = get\_lock(B)}   & \texttt{lock2 = get\_lock(A)} \\
$t_3$ & \texttt{acquire(lock1)} $\to$ obtiene \emph{lock\_A} & \texttt{acquire(lock1)} $\to$ obtiene \emph{lock\_B} \\
$t_4$ & \texttt{acquire(lock2)} (pide lock\_B) & \texttt{acquire(lock2)} (pide lock\_A) \\
$t_5$ & \textbf{BLOQUEA}: lock\_B lo tiene T2 & \textbf{BLOQUEA}: lock\_A lo tiene T1 \\
\bottomrule
\end{tabular}
\end{center}

Resultado: ambos threads quedan bloqueados, esperando un lock que
el otro retiene. Se cumplen las cuatro condiciones de Coffman:

\begin{itemize}[leftmargin=*,nosep]
    \item \textbf{Exclusión mutua}: cada mutex sólo puede pertenecer
        a un thread a la vez.
    \item \textbf{Hold and wait}: cada thread tiene un lock y espera
        otro.
    \item \textbf{No preemption}: el sistema no puede quitarle el
        mutex a un thread.
    \item \textbf{Espera circular}:
        T1 espera lock\_B (de T2) $\to$ T2 espera lock\_A (de T1)
        $\to$ T1.
\end{itemize}

\subsection*{Grafo de espera}

\begin{center}
\begin{tikzpicture}[
    every node/.style={font=\ttfamily\small},
    proc/.style={circle,draw,minimum size=0.9cm,font=\bfseries\small},
    res/.style={rectangle,draw,minimum size=0.9cm,font=\bfseries\small},
    every edge/.style={draw,-{Stealth[length=2.5mm]},thick}
]
\node[proc] (T1) at (0, 2)    {T1};
\node[proc] (T2) at (4, 2)    {T2};
\node[res]  (LA) at (-1.2, 0) {lock\_A};
\node[res]  (LB) at (5.2, 0)  {lock\_B};

\draw[->,blue!60!black,thick]   (LA) -- node[left,font=\scriptsize]{tiene} (T1);
\draw[->,blue!60!black,thick]   (LB) -- node[right,font=\scriptsize]{tiene} (T2);
\draw[->,red!70!black,thick,dashed]   (T1) -- node[above,font=\scriptsize]{espera} (LB);
\draw[->,red!70!black,thick,dashed]   (T2) -- node[above,font=\scriptsize]{espera} (LA);
\end{tikzpicture}
\end{center}

\subsection*{Cómo prevenirlo}

La técnica estándar es imponer un \textbf{orden global de
adquisición} de los locks, independiente del orden de los argumentos.
Una forma simple es comparar por identidad (por ejemplo, dirección
de memoria o ID numérico de la cuenta) y tomar primero el menor:

\begin{lstlisting}[style=cstyle]
void transfer(Account from, Account to, double amount)
{
    mutex first, second;

    /* Orden global: tomar primero el mutex de la cuenta con menor id */
    if (id(from) < id(to)) {
        first  = get_lock(from);
        second = get_lock(to);
    } else {
        first  = get_lock(to);
        second = get_lock(from);
    }

    acquire(first);
    acquire(second);

    withdraw(from, amount);
    deposit(to, amount);

    release(second);
    release(first);
}
\end{lstlisting}

Con esta corrección, \emph{ambos} threads ---no importa quién haga
A$\to$B ni B$\to$A--- tomarían \emph{primero} lock\_A y \emph{después}
lock\_B. Es imposible que se forme un ciclo.

\begin{respuesta}
\textbf{NO previene deadlocks.} El orden de adquisición de los
mutex depende de cuál cuenta sea \texttt{from} y cuál \texttt{to};
si dos hilos transfieren en sentidos contrarios
(\texttt{A}$\to$\texttt{B} y \texttt{B}$\to$\texttt{A})
simultáneamente, pueden tomar los locks en orden inverso y caer en
una espera circular clásica.

\textbf{Solución}: imponer un \emph{orden global} de adquisición de
los locks (por ejemplo, por ID de cuenta ascendente), independiente
de cuál cuenta sea \texttt{from} o \texttt{to}.
\end{respuesta}

% =====================================================================
\section{Ejercicio 5 --- Round Robin con quantum = 2 (2 pts)}
% =====================================================================

\textbf{Enunciado.} Asumiendo que los procesos están encolados en
orden P1, P2, P3, dar un diagrama de Gantt de la planificación según
un planificador Round Robin (RR) con \texttt{quantum = 2}.

\begin{center}
\small
\begin{tabular}{@{}cc@{}}
\toprule
Proceso & Ráfaga de uso de CPU \\
\midrule
P1 & 5 \\
P2 & 2 \\
P3 & 3 \\
\bottomrule
\end{tabular}
\end{center}

\subsection*{Análisis: cómo planifica RR}

Round Robin mantiene una cola FIFO de procesos \texttt{READY}. El
scheduler asigna al primero un \emph{quantum} de tiempo de CPU; al
agotarse:

\begin{itemize}[leftmargin=*,nosep]
    \item Si el proceso tiene más ráfaga restante, se lo
        \emph{desaloja} (\emph{preempt}) y se lo encola al
        \textbf{final} de la lista.
    \item Si la ráfaga se completó antes o exactamente al quantum,
        el proceso \textbf{termina} (en este ejercicio no hay I/O).
\end{itemize}

Todos los procesos arriban en $t = 0$ y entran en el orden P1, P2,
P3 a la cola \texttt{READY}.

\subsection*{Traza paso a paso}

Notación: \texttt{[P:rem]} indica el proceso \texttt{P} con
\texttt{rem} unidades de ráfaga \emph{remanentes}.

\begin{center}
\small
\begin{tabular}{@{}clp{4cm}l@{}}
\toprule
Intervalo & CPU & Evento al final del intervalo & Cola \texttt{READY} \\
\midrule
$[0, 2)$  & P1 & cuantum agotado, queda $5-2=3$ & \texttt{[P2:2, P3:3, P1:3]} \\
$[2, 4)$  & P2 & P2 termina exactamente (ráfaga $=2$) & \texttt{[P3:3, P1:3]} \\
$[4, 6)$  & P3 & cuantum agotado, queda $3-2=1$ & \texttt{[P1:3, P3:1]} \\
$[6, 8)$  & P1 & cuantum agotado, queda $3-2=1$ & \texttt{[P3:1, P1:1]} \\
$[8, 9)$  & P3 & ráfaga termina (queda $0$), P3 \textbf{EXIT}  & \texttt{[P1:1]} \\
$[9, 10)$ & P1 & ráfaga termina (queda $0$), P1 \textbf{EXIT}  & \texttt{[\ ]} \\
\bottomrule
\end{tabular}
\end{center}

Verificación: la CPU está ocupada de $0$ a $10$ sin huecos, y el
total de CPU usada es $5 + 2 + 3 = 10$. ✓

\subsection*{Diagrama de Gantt}

\begin{center}
\begin{tikzpicture}[xscale=1.0,yscale=0.7,every node/.style={font=\ttfamily\small}]
\draw[fill=c1!50] (0,0)  rectangle (2,1)  node[midway,white,font=\bfseries]{P1};
\draw[fill=c2!50] (2,0)  rectangle (4,1)  node[midway,white,font=\bfseries]{P2};
\draw[fill=c3!50] (4,0)  rectangle (6,1)  node[midway,white,font=\bfseries]{P3};
\draw[fill=c1!50] (6,0)  rectangle (8,1)  node[midway,white,font=\bfseries]{P1};
\draw[fill=c3!50] (8,0)  rectangle (9,1)  node[midway,white,font=\bfseries\small]{P3};
\draw[fill=c1!50] (9,0)  rectangle (10,1) node[midway,white,font=\bfseries\small]{P1};

\foreach \x in {0,2,4,6,8,9,10}
    \draw (\x,0) -- (\x,-0.2) node[below,font=\small]{\x};
\end{tikzpicture}
\end{center}

\subsection*{Métricas}

Aplicando turnaround $=$ Fin $-$ Llegada y espera $=$ turnaround $-$
ráfaga:

\begin{center}
\small
\begin{tabular}{@{}lccccc@{}}
\toprule
Proceso & Llegada & Fin & Turnaround & Ráfaga & Espera \\
\midrule
P1 & 0 & 10 & 10 & 5 & 5 \\
P2 & 0 & 4  & 4  & 2 & 2 \\
P3 & 0 & 9  & 9  & 3 & 6 \\
\bottomrule
\end{tabular}
\end{center}

\textbf{Promedios:}

\quad Turnaround promedio = $(10 + 4 + 9) \div 3 = 23 \div 3 ≈ 7.67$

\quad Espera promedio = $(5 + 2 + 6) \div 3 = 13 \div 3 ≈ 4.33$

\begin{respuesta}
Gantt: \texttt{P1 | P2 | P3 | P1 | P3 | P1} con marcas de tiempo
$0, 2, 4, 6, 8, 9, 10$.

Tiempos de finalización: $\textbf{P2 = 4}$, $\textbf{P3 = 9}$,
$\textbf{P1 = 10}$. P2 (la ráfaga más corta) sale primero pese a no
ser la primera en encolarse; P1 (la más larga) sale al final.

Promedios: turnaround $≈ 7.67$, espera $≈ 4.33$.
\end{respuesta}

\begin{idea}
RR es \textbf{justo} (\emph{fair}): nadie monopoliza la CPU porque
los desalojos son periódicos. Pero \emph{no} optimiza el turnaround
promedio: SJF haría P2 antes (orden P2, P3, P1) y daría mejor
promedio. RR sacrifica eficiencia por baja latencia de respuesta,
ideal para sistemas \emph{time-sharing} interactivos.
\end{idea}

% =====================================================================
\section*{Resumen del recuperatorio}
\addcontentsline{toc}{section}{Resumen del recuperatorio}
% =====================================================================

\begin{center}
\small
\begin{tabular}{@{}clp{8.5cm}@{}}
\toprule
\# & Tema & Respuesta sintética \\
\midrule
1 & V/F syscalls         & A:F (\texttt{exec} no crea), B:F (CRT antes de main), C:F (\texttt{printf} es libc), D:V, E:F (COW + redacción), F:V \\
2 & Conteo de \texttt{fork} & $2^3 = 8$ procesos totales $\Rightarrow$ \textbf{7 nuevos} (sin contar el original) \\
3 & Pipe bidireccional   & Incorrecto: un solo pipe no separa sentidos. Corrección: dos pipes, uno por sentido \\
4 & Deadlock en \texttt{transfer} & No previene; escenario T1: A$\to$B y T2: B$\to$A. Fix: orden global de locks por ID \\
5 & Round Robin q=2      & Gantt: P1,P2,P3,P1,P3,P1; fin P2=4, P3=9, P1=10; turnaround promedio $\approx 7.67$ \\
\bottomrule
\end{tabular}
\end{center}

\end{document}
